\documentclass[10pt]{article}
\usepackage[margin=0.72in]{geometry}
\usepackage[utf8]{inputenc}
\usepackage[T1]{fontenc}
\usepackage{hyperref}
\usepackage{url}
\usepackage{booktabs}
\usepackage{amsfonts}
\usepackage{microtype}
\usepackage{xcolor}
\usepackage{graphicx}
\usepackage{float}
\usepackage{amsmath}
\usepackage{subcaption}
\usepackage{placeins}
\raggedbottom
\graphicspath{{./}{figures/}}
\setlength{\textfloatsep}{5pt plus 1pt minus 1pt}
\setlength{\intextsep}{5pt plus 1pt minus 1pt}
\setlength{\floatsep}{5pt plus 1pt minus 1pt}
\setlength{\abovecaptionskip}{2pt}
\setlength{\belowcaptionskip}{0pt}
\usepackage{titlesec}
\titlespacing*{\section}{0pt}{6pt}{2pt}
\titlespacing*{\subsection}{0pt}{4pt}{1pt}
\definecolor{good}{HTML}{1F8A5A}
\definecolor{bad}{HTML}{C0392B}
\newcommand{\yes}{\textcolor{good}{\textbf{helps}}}
\newcommand{\no}{\textcolor{bad}{\textbf{fails}}}

\title{Adversarial Co-Evolution of RL and LLM Agents in Gin Rummy:\\
       A Gold-Standard Study of What Makes a Lightweight Agent Strong}

\author{Nima Kelidari \quad Mahdi Salmani \quad Mohammadsaeed Haghi \\
  \texttt{\{kelidari, salmanis, haghim\}@usc.edu} \\
  University of Southern California}
\date{}

\begin{document}
\maketitle

\begin{abstract}
Training a reinforcement-learning (RL) agent for a card game with hidden information runs into the
\emph{opponent bottleneck}: an agent is only as good as the opponents it practises against. We study
this in Gin Rummy, used here as a clean example and not as the goal. We build a full framework with
five parts: an action-masked PPO/TRPO learner, a hand-coded \emph{gold-standard} (perfect) opponent
used only as a measuring stick, a distributed server that lets a 7B language model (LLM) play inside
the game, an opponent-curriculum system, and a browser game. We then test almost every reasonable
way to make the agent stronger, judged on one number: win-rate against the perfect player. Across
$100+$ controlled runs we find three things. First, the gold standard beats every learned agent
($70$ to $99\%$) yet gins in under $2\%$ of games, winning by knocking with low deadwood, and an RL
agent learns the same habit even when a gin pays $3\times$ a knock. Second, imitation learning
(DAgger), dense short-term reward shaping, and training against a live LLM all fail to beat plain
curriculum RL, each for a clear reason. Third, the ideas that \emph{do} help, namely a trust-region
algorithm (TRPO), a knock-first reward, a rising opponent curriculum, warm-starting, and keeping the
best checkpoint, stack together and push a small agent from the $\sim\!30\%$ self-play champion to
$\mathbf{34\%}$ against perfect play. We report the engineering and an honest look at why the popular
hybrid methods do not beat pure RL here. The framework and these lessons are not specific to Gin
Rummy. They carry over to other two-player games and to RL plus LLM agent systems in general.
\end{abstract}

\section{Introduction}
Reinforcement learning is good at making a sequence of decisions, but it needs a lot of data and it
depends heavily on its opponents. Train against a weak player and the agent learns to beat weakness
rather than to play well. Strong opponents are exactly what is hard to get: hand-built solvers are
specific to one game, and humans do not scale. Large language models (LLMs) look like a shortcut
because they bring broad, common-sense strategy with no training. This project asks, in the simple
testbed of Gin Rummy, \textbf{how close a small, fast RL agent can get to perfect play, and which
training ideas actually move it there.} We pick Gin Rummy because it is easy to score yet hard to
play well, but the questions and the answers are general.

\noindent\textbf{Contributions.} (1) A reusable framework that ties together masked-PPO/TRPO RL, a
\emph{perfect} benchmark agent, a distributed LLM-serving stack, an opponent curriculum, and a web
game (Sec.~2). (2) A gold-standard benchmark and one clean result: the perfect player almost never
gins, and an RL agent learns the same thing no matter the reward (Sec.~3, 4.6). (3) A side-by-side
study of about ten training methods on one number, win-rate against the perfect player, with the
\emph{reason} each one works or fails (Sec.~4, Fig.~\ref{fig:regimes}), including honest negatives
for DAgger, dense rewards, learned embeddings, and a live LLM in the loop. (4) A recipe that stacks
the helpful ideas to reach $34\%$ against perfect play (Sec.~4.7, 5).

\section{The Framework}
\textbf{Environment and learner.} We use PettingZoo/RLCard \texttt{gin\_rummy\_v4}. The observation
is four $52$-card planes (hand, top discard, known picks, the rest). Our \texttt{MaskedGinRummyPolicy}
sets illegal-action scores to $-\infty$, so the agent only ever picks a legal move, and it acts by
choosing the highest-probability \emph{legal} move (no random fallback). \textbf{Gold standard.} A
hand-coded perfect player works out the exact best melds with RLCard's solver and knocks the moment
it should. It is the measuring stick and \emph{never} trains the RL agent. \textbf{Distributed LLM.}
One rollout asks the opponent tens of thousands of questions, and at $0.5$ to $3$\,s per 7B call a
simple loop takes hours, so we keep serving separate from training. A master node balances load over
self-registering GPU workers behind a suit-symmetry cache, using the normal Ollama API. One
infrastructure result stood out: putting 7B weights on scratch (BeeGFS) instead of home storage cuts
worker load from about $28$\,min to about $27$\,s, a $62\times$ speedup, and a $14$-worker pool
handled about $32$ questions per second. \textbf{Curriculum and game.} A scheduler ramps opponents
(random, then past selves, then strong models), and a zero-dependency browser client lets people play
any agent. Fig.~\ref{fig:arch} shows how the parts connect.

\begin{figure}[h]\centering
\includegraphics[width=0.80\linewidth]{architecture.png}
\caption{System overview. The learner trains in the environment against a curriculum of opponents
(and, in dedicated runs, the LLM); the gold standard is used for \emph{scoring only}.}\label{fig:arch}
\end{figure}

\section{The Gold-Standard Benchmark and a Surprise}
To grade everything on one fair scale we built a perfect player and ran it against the whole set of
trained agents (Fig.~\ref{fig:gold}, $600$ games each). It wins $70.2\%$ against the self-play
champion and $70$ to $99\%$ against the rest. The surprise is how it wins: the perfect player
\textbf{gins only $0.7$ to $1.7\%$ of the time}, even though gin scores the most points. It wins by
\emph{knocking early with low deadwood}. That goes against the natural guess that chasing gin is
good, and it changes how we read every reward experiment below.

\begin{figure}[h]\centering
\includegraphics[width=0.84\linewidth]{gold_bench.png}
\caption{The perfect player beats every learned agent (left) yet gins under $2\%$ of the time
(right): it wins by low-deadwood knocking, not by chasing gin.}\label{fig:gold}
\end{figure}

\section{What We Tried: A Side-by-Side Study}
We tested almost every reasonable training method on the same number, best win-rate against the
perfect player (Fig.~\ref{fig:regimes}), and for each we give the reason it lands where it does.

\begin{figure}[h]\centering
\includegraphics[width=0.76\linewidth]{regimes.png}
\caption{Every method we tested, ranked by measured win-rate against the perfect player.}
\label{fig:regimes}
\end{figure}

\subsection{RL backbone and reward shaping}
Masked PPO maxes out against random opponents: an eight-setting sweep lands every run in the $98.3$ to
$99.6\%$ band (you cannot tell them apart at $1000$-game eval), and a larger set of $40+$
reward-shaping runs shows the gin-to-knock reward ratio is a real dial. A knock-first reward gives
$97\%$ knocks (win $99.4\%$), while a gin-first reward gives $22\%$ gins (win $98.3\%$). \emph{Reason:}
random play is too weak to teach strategy, so the easy number maxes out, and the reward ratio sets how
much risk the agent takes. This is why we need stronger opponents and a sensible reward.

\subsection{Self-play and curriculum}
Self-play is a free, rising curriculum: a $3$M-step self-play agent beats its own parent $61\%$ of the
time. But a pool of past selves with no guidance \emph{broke down} after about $10$M steps.
\emph{Reason:} self-play makes difficulty for free, but it can chase circular strategies (where A
beats B, B beats C, and C beats A), so the schedule and the mix of opponents matter. That motivates the controlled sweep in
Sec.~4.6.

\subsection{Why the hybrid methods fail (honest negatives)}
\textbf{Live LLM in the loop.} Mid-size LLMs (Qwen2.5-7B, gpt-oss-20b) play real Gin Rummy with a
chain-of-thought prompt ($100\%$ legal moves) and even beat our self-play agent $3$ to $2$ in short
matches. But at $9$ to $27$\,s per move they are far too slow for the millions of moves RL needs, and
even with the distributed stack, training this way would take weeks. Vision-LLMs failed outright
(about $10\%$ legal moves): this is a text task. \textbf{Imitation (DAgger).} Copying an expert's
moves one at a time collapses to almost no wins. \emph{Reason:} the student copies moves in familiar
spots but never learns the thinking behind them (like tracking the opponent), so it falls apart in new
situations. Copying a move is not the same as understanding why. \textbf{Dense rewards.} Rewarding
every step (horizon $L\!=\!2,5$) makes the agent short-sighted: $L\!=\!2$ shows no learning and
$L\!=\!5$ stops improving after about $500$k steps, farming instant points instead of trying to win.
\emph{Reason:} dense shaping pulls the agent away from the real, sparse goal. The rule that follows:
we let the LLM or expert \emph{inform our design}, but never supervise the agent directly.

\subsection{Algorithm: PPO vs.\ TRPO}
With the same masking, curriculum, and reward, TRPO beats PPO at $2$M steps on every opponent and both
seeds: $50.5\%$ vs.\ $31.0\%$ against the champion and $22.5\%$ vs.\ $15.0\%$ against gold
(Fig.~\ref{fig:algo}a). \emph{Reason:} TRPO takes smaller, safer steps when it updates the policy,
which fits a setting where rewards are rare and the opponent keeps changing. (GRPO and DPO do not
apply here: they align language models and have no per-move game version.)

\subsection{State representation}
We squeezed the big, sparse observation into a small dense $20$-d vector two ways: one that learned
which game states are close together, and one where an LLM judged how similar two states are. Used
frozen as the agent's input, both \emph{lose} to the raw sparse input against gold ($6$ to $14\%$ vs.\
$15\%$). \emph{Reason:} a fixed, squeezed vector that is good at ``are these two states similar'' still
throws away detail the agent needs, and neither a bigger size nor un-freezing fixed it.

\begin{figure}[h]\centering
\begin{subfigure}{0.49\linewidth}\centering\includegraphics[width=\linewidth]{algo_compare.png}
\caption{TRPO beats PPO vs.\ every opponent.}\end{subfigure}\hfill
\begin{subfigure}{0.49\linewidth}\centering\includegraphics[width=\linewidth]{reward_gin.png}
\caption{Gin-rate stays $<\!1\%$ for every reward.}\end{subfigure}
\caption{Algorithm and reward studies (2 seeds each).}\label{fig:algo}
\end{figure}

\subsection{The curriculum sweep and the key finding}
A $30$-run sweep that changes one thing at a time across algorithm, reward, and curriculum lands every
recipe near champion strength ($\sim\!30$ to $33\%$ vs.\ gold), but it gives one clean, solid result
(Fig.~\ref{fig:algo}b): \textbf{the gin-rate stays under $1\%$ no matter the reward}. Even paying
$3\times$ for a gin does not raise it above the balanced baseline. \emph{Reason:} like the perfect
player, the agent works out that chasing gin loses, and you cannot pay it into the bad habit. The
reward that does help is the opposite, a small ``knock earlier'' penalty that gives the shortest games
and the best play. Fig.~\ref{fig:learn} shows the mechanism behind the numbers: win-rate against the
champion climbs as the curriculum swaps in tougher opponents (random, then pool, then self, then
strong), and the late dip on one run is the drift that keep-the-best fixes (Sec.~4.7).

\begin{figure}[h]\centering
\includegraphics[width=0.64\linewidth]{learning_curves.png}
\caption{Skill (win-rate vs.\ the champion) rises through the curriculum stages; the late dip
motivates best-checkpoint selection.}\label{fig:learn}
\end{figure}

\subsection{Keep-the-best and warm-start (the recipe)}
A final series stacks the ideas that help: \emph{warm-start} from the champion, the knock-first reward
plus a small \emph{deadwood-reduction} reward (a nudge toward the optimal style that does not change
which strategy is best, computed from the agent's own public hand), a curriculum seeded with the
strongest earlier agents, and, most important, \emph{keeping the best checkpoint} (training drifts past
its peak, so we evaluate every $1$M steps and ship the best one). Checked over $2000$ games, the
strongest agent reaches $\mathbf{34.2\% \pm 2.1\%}$ vs.\ gold and $51.4\%$ vs.\ the old champion, our
best, up from the champion's $\sim\!30\%$.

\section{What Mattered, and the Bottom Line}
\begin{table}[h]\centering\footnotesize
\begin{tabular}{@{}llp{8.2cm}@{}}\toprule
Idea & Verdict & Why \\\midrule
Keep the best checkpoint & \yes & training drifts past its peak; saving the best recovers $2$ to $3$ pts \\
Warm-start from the champion & \yes & start strong, then improve \\
TRPO over PPO & \yes & smaller, safer steps suit rare rewards and a changing opponent \\
Reward knocking, not gin & \yes & copies the perfect player's low-risk style \\
Rising opponent curriculum & \yes & always a fair but harder challenge \\
Pay more for gin & \textcolor{bad}{no effect} & the agent refuses the bad habit at any reward \\
Learned state embeddings & \no & a fixed, squeezed input throws away useful detail \\
Imitation (DAgger) & \no & copies moves, not the thinking behind them \\
Dense short-term rewards & \no & short-sighted: greedy for points, blind to winning \\
Live LLM in the loop & \no & strong, but far too slow for millions of moves \\
\bottomrule\end{tabular}
\caption{Which ideas moved the needle against the perfect player, across $100+$ runs.}
\label{tab:levers}
\end{table}

Stacking the helpful ideas took a small agent from the champion's $\sim\!30\%$ to $34\%$ against
\emph{perfect} play (Fig.~\ref{fig:journey}) and gave one reusable result (the best play rarely gins).
The honest ceiling is the mid-$30$s percent: tuning reward, algorithm, and opponents tops out there.
None of this is specific to Gin Rummy. The same recipe, a strong reference to grade against, a rising
ladder of opponents, a reward that copies strong play, and best-checkpoint selection, applies to other
two-player games and to RL plus LLM agent systems in general.

\begin{figure}[h]\centering
\includegraphics[width=0.54\linewidth]{journey.png}
\caption{Closing the gap to perfect play, step by step.}\label{fig:journey}
\end{figure}

\section{Conclusion and Future Work}
We built the framework, the benchmark, and the side-by-side study to answer, fairly, how strong a
small RL agent can be in Gin Rummy and what makes it so. Plain curriculum RL, sharpened with a
trust-region algorithm, a knock-first reward, warm-starting, and best-checkpoint selection, is the
winner, and the popular hybrid shortcuts (imitation, dense reward, a live LLM in the loop) each fail
for a clear, named reason. Going clearly past the mid-$30$s ceiling against perfect play almost
certainly needs a \emph{different kind} of method, such as search and planning (the kind that beat the
best humans at poker, counterfactual-regret minimization) or an agent that remembers the whole game so
far, rather than more reward tuning. The framework, the benchmark, and the distributed LLM data
pipeline are the reusable base for that next step (for example, offline RL on a large LLM-vs-LLM
dataset), and they apply well beyond this one game. All figures rebuild from saved JSON, and live
training curves are on Weights \& Biases.

\begin{thebibliography}{9}\footnotesize
\bibitem{ppo} J.~Schulman et al. Proximal Policy Optimization Algorithms. \emph{arXiv:1707.06347}, 2017.
\bibitem{trpo} J.~Schulman et al. Trust Region Policy Optimization. \emph{ICML}, 2015.
\bibitem{mask} S.~Huang and S.~Onta\~{n}\'{o}n. A Closer Look at Invalid Action Masking in Policy Gradient Algorithms. \emph{arXiv:2006.14171}, 2022.
\bibitem{dagger} S.~Ross, G.~Gordon, D.~Bagnell. A Reduction of Imitation Learning to No-Regret Online Learning (DAgger). \emph{AISTATS}, 2011.
\bibitem{cfr} N.~Brown and T.~Sandholm. Superhuman AI for multiplayer poker. \emph{Science} 365(6456):885--890, 2019.
\bibitem{pz} J.~Terry et al. PettingZoo: Gym for Multi-Agent Reinforcement Learning. \emph{NeurIPS}, 2021.
\bibitem{rlcard} D.~Zha et al. RLCard: A Toolkit for Reinforcement Learning in Card Games. \emph{arXiv:1910.04376}, 2019.
\bibitem{sb3} A.~Raffin et al. Stable-Baselines3. \emph{JMLR} 22(268):1--8, 2021.
\bibitem{qwen} Qwen Team. Qwen2.5 Technical Report. \emph{arXiv:2412.15115}, 2024.
\end{thebibliography}

\end{document}
